\section*{Kapitel 6 --- Interfacet och stubben}
\addcontentsline{toc}{section}{Kapitel 6 --- Interfacet och stubben}

\solution{ex:6:1}{Förutsäg kompilatorn}
\ans{a} Det kompilerar utan minsta invändning, och det är hela problemet. Utan \code{virtual}
bestäms vilken destruktor som anropas av pekarens \emph{statiska} typ, alltså \code{Interface}.
\code{Stub}s destruktor körs aldrig, dess medlemmar destrueras aldrig, och standarden kallar det
odefinierat beteende. För \code{Stub} märks det inte --- den äger ingenting --- men för en klass
som håller en resurs läcker den, tyst, vid varje radering. Inget test fångar det.
\ans{b} Utan attributet är \code{driver.send(frame);} en fullt laglig sats som kompilatorn inte
har någon anledning att klaga på. Med \code{[[nodiscard]]} blir den en varning, och med
\code{-Werror} ett fel. Det förvandlar den vanligaste drivrutinsbuggen --- att tro att en frame
gick iväg när hårdvaran var upptagen --- från en felsökningssession till en kompilering som
stannar.
\ans{c} Med \code{override} är felstavningen ett kompileringsfel: det finns ingen virtuell metod
med det namnet att överskugga. Utan \code{override} kompilerar det, och du har i stället
\emph{lagt till} en ny metod. Basklassens rena virtuella metod är då fortfarande oimplementerad,
så klassen förblir abstrakt --- och felmeddelandet du får handlar om att typen inte går att
instansiera, långt ifrån den felstavade raden. Det är därför \code{override} är obligatoriskt.

\solution{ex:6:2}{En metod för mycket}
\ans{a} Ingenting meningsfullt. Stubben har inga register; den skulle behöva returnera en påhittad
konstant, alltså ljuga.
\ans{b} \code{EchoNode} blir bunden till att hårdvaran \emph{har} register. Testet på om en metod
hör hemma i interfacet är just detta: kan båda implementationerna uppfylla den ärligt? Kan de inte
det, hör den inte hemma där.
\ans{c} Varje anropare måste ändras, eftersom bitarnas betydelse är en del av registerkartan och
inte av interfacet. Interfacet skulle alltså ärva hårdvarans versioner, vilket är precis vad det
finns för att undvika.
\ans{d} I \code{Spi}, som privat hjälpfunktion --- eller, för felsökning, i en separat
\code{PrintTransport} under \code{ByteTransport}-sömmen, där den inte syns för någon ovanför.

\solution{ex:6:3}{Polymorfismen}
Uppgiften ber inte om kod; frågorna besvaras här.
\ans{a} Nej. Funktionen tar en \code{driver::can::Interface\&} och använder bara interfacets fyra
metoder.
\ans{b} Därför att den då kan köras mot vilken implementation som helst. Skriven mot \code{Stub}
hade den varit ett test av \code{Stub}; skriven mot interfacet är den ett test av
\emph{kontraktet}, och kontraktet är det både \code{Stub} och \code{Spi} lovar.
\ans{c} Att \code{Spi} verkligen uppfyller samma kontrakt: samma returvärden i samma situationer.
Framför allt måste \code{send()} returnera \code{false} när hårdvaran är upptagen och
\code{receive()} returnera \code{false} när ingen frame finns --- annars går funktionen igenom för
stubben och fallerar för drivrutinen.

\solution{ex:6:4}{Stubbens gränser}
\ans{a} Fel (\code{hasError()} är alltid \code{false}), upptagen hårdvara (\code{send()} lyckas
alltid), kö (en andra \code{send()} skriver över den första), och allt CAN gör: arbitrering, CRC,
bitstoppning.
\ans{b} Till exempel upptagen hårdvara: en räknare som får nästa \code{send()} att returnera
\code{false}, satt av en testmetod. Det är den enklaste, eftersom den inte kräver något nytt
tillstånd utöver en flagga.
\ans{c} Varje simulerat beteende är kod som kan vara fel. En stubb som simulerar kö, fel och
timing är en andra implementation av hårdvaran --- och då vet du inte längre om ett fallerande test
beror på drivrutinen eller på stubben. Gränsen går ungefär där stubben skulle behöva ett eget
testfall: behöver den det, hör beteendet hemma i en riktigare dubblare, som
\code{BankTransport} i \kapref{ch:testning}.
